\section{Introduction}
\label{sec:introduction}

Generative models have made remarkable progress on pixel-space images, but the
computational cost of operating directly on dense pixel grids motivates the search
for compact latent representations.
\emph{Gaussian splatting}~\cite{kerbl20233d} provides one such representation:
an image is described as a superposition of $K$ anisotropic 2-D Gaussians,
each parameterised by position, scale, orientation, colour, and opacity.
Such representations are interpretable, differentiable, and naturally sparse.

The \textbf{GaussianDiffusion} project exploits Gaussian splatting as a latent
space for a denoising diffusion probabilistic model (DDPM).  The pipeline is:

\begin{enumerate}[leftmargin=2em]
  \item \textbf{Encode:} fit $K$ Gaussians to each training image using gradient
        descent $\rightarrow$ tensor $\W \in \R^{K \times 7}$.
  \item \textbf{Train:} learn a DDPM GaussianTransformer that generates $\W$
        tensors from Gaussian noise.
  \item \textbf{Sample:} run reverse diffusion on the transformer, then render
        the predicted $\W$ with the Gaussian splatting renderer to obtain images.
\end{enumerate}

\noindent
This report focuses exclusively on \textbf{Step~1}: the correctness, quality, and
downstream suitability of the encoder.
A sound encoder is a prerequisite for everything downstream.
Poor reconstructions set a hard ceiling on generative quality; parameter
distributions that violate assumptions (non-smooth, multi-modal, poorly normalised)
make diffusion training harder.

\subsection{Scope and Contributions}

\noindent This report makes the following contributions:

\begin{itemize}[leftmargin=2em]
  \item A full description of the encoder's architecture and optimisation procedure
        (\S\ref{sec:architecture}).
  \item A mathematical account of the Gaussian splatting renderer, including an
        explanation of the \texttt{F.affine\_grid} coordinate convention
        (\S\ref{sec:renderer}).
  \item Discovery and diagnosis of a coordinate-inversion bug that prevented
        convergence, and verification of the three-line fix (\S\ref{sec:bug}).
  \item Quantitative validation on MNIST: per-digit PSNR statistics, convergence
        dynamics, K-sensitivity, loss-function ablation, and parameter analysis
        (\S\ref{sec:experiments}).
  \item Assessment of downstream readiness for diffusion model training
        (\S\ref{sec:downstream}).
\end{itemize}

\subsection{Dataset and Setting}

All experiments use the MNIST handwritten-digit dataset
($28 \times 28$ greyscale images, pixel values in $[0,1]$).
We encode $N = 50$ images (5 per digit class) using the corrected encoder
with $K = 70$ Gaussians, $3{,}000$ maximum epochs, learning rate $5\!\times\!10^{-3}$,
kernel size $11$, and early stopping at MSE\,$ < 10^{-4}$ ($\approx 40\,\text{dB}$
PSNR).  Additional experiments vary $K$, the loss function, and the random seed.
All computations are performed on CPU; the same code runs on GPU without modification.
